\centerline{\bf \Huge Эйлер и Рыцарь}


\begin{flushright}
        \scriptsize{Давайте решим о будущем в будущем}
\end{flushright}

\problem{0.1}
Давным давно островитянин Дерб сказал своим друзьям: - Вчера мой сосед заявил мне, что он лжец! Кем является Дерб — рыцарем или лжецом?
% Лжец

\problem{0.2}
Каждый из 7 сидящих за круглым столом жителей острова сказал: 'Мои соседи лжец и рыцарь'. Сколько рыцарей и сколько лжецов сидит за столом?
% 2

\problem{0.3}
У Пети есть карточки, у которых каждая сторона — красная или синяя. Карточек с синей стороной — 10, карточек с красной стороной — 12, карточек с разными сторонами — 5. Сколько карточек, у которых стороны одного цвета?

\problem{1}
Однажды в четверг после дождя между островитянами Тимом и Томом произошел следующий диалог: - Ты можешь сказать, что я рыцарь, - гордо заявил Тим. - Ты можешь сказать, что я лжец, - грустно ответил ему Том. Кем являются Тим и Том?
% Тим - лжец 
% Том - рыцарь

\problem{2}
За столом по кругу сидит 10 аборигенов Каждый из них произнес фразу: «Следующие 4 человека, стоящие после меня по часовой стрелке, лжецы». Сколько среди них лжецов?
% 2

\problem{3}
Каждый из 9 сидящих за круглым столом жителей острова сказал: 'Мои соседи лжец и рыцарь'. Сколько рыцарей и сколько лжецов сидит за столом?
%  0 рыцарей, 9 лжецов ИЛИ 6 рыцарей, 3 лжеца.

\problem{4}
В классе все увлекаются математикой или биологией. Сколько человек в классе, если математикой увлекаются 15 человек, биологией — 20, а математикой и биологией одновременно — 10?


\problem{5}
В классе 33 ученика. 27 из них изучают английский язык, и 10 — французский. Сколько учеников могут изучать и английский, и французский?



\problem{6}
Среди монстров, живущих в подземелье, у 14 есть крылья, у 13 — рога, у 15 — хвост, у 6 — крылья и рога, у 5 — рога и хвост, у 7 — крылья и хвост. Кроме того, у двух монстров есть и крылья, и рога, и хвост, и ещё у двух нет ни крыльев, ни рогов, ни хвоста. Сколько всего монстров живёт в подземелье?

\problem{ы}
Каждый из трёх игроков записывает 100 слов, после чего записи срав- нивают. Если слово встретилось хотя бы у двоих, то его вычеркивают из всех списков. Могло ли так оказаться, что у первого игрока осталось 54 сло- ва, у второго — 75 слов, а у третьего — 80 слов?
